\begin{statement}{C.2.UN.23}{Proved}{A cyclic torus quotient with nonnegative Kodaira dimension is a torus, bielliptic, or K3}
Let $A$ be a compact complex two-torus, let $G$ be a finite
cyclic group of holomorphic automorphisms acting effectively
on $A$, and let $Z\to A/G$ be a smooth resolution.
Every positive plurigenus of $Z$ is at most one. There are
the following exhaustive alternatives for a generator
$g(z)=Lz+b$:
\[
\begin{array}{c|c}
\text{linear part and action}&\text{minimal resolution type}\\ \hline
L=I&\text{a two-torus}\\
\dim_{\C}\ker(L-I)=1,\ G\text{ acts freely}
   &\text{a bielliptic surface}\\
\ker(L-I)=0,\ \det L=1&\text{a K3 surface}\\
\text{all other cases}&P_m(Z)=0\text{ for every }m>0 .
\end{array}                                                    \tag{C.CQ.1}
\]
In particular an Enriques surface is not bimeromorphic to
a cyclic quotient of a complex two-torus.

For the actual component cover in C.2.UN.16, suppose its
minimal model is a two-torus. The quotient by its cyclic
stabilizer, followed by the specified proper birational
comparison to the original normalization, satisfies
(C.CQ.1). Thus an original component of nonnegative
Kodaira dimension in this branch has minimal model a
torus, a bielliptic surface, or a K3 surface.
\end{statement}
\begin{proof}
The derivative of a holomorphic torus automorphism is
constant in the translation trivialization: its entries
are holomorphic functions on a compact torus. Lifting to
the universal cover consequently writes the automorphism
as $z\mapsto Lz+b$, with $L\Lambda=\Lambda$ for the
lattice $\Lambda$. For a finite-order automorphism, $L$
is diagonalizable with roots of unity as eigenvalues.

We first specify the canonical comparison used in every
case. Resolve the rational lift from $A$ to $Z$:
\[
\begin{array}{ccc}
A'&\xrightarrow{\ v\ }&Z\\
\downarrow{\scriptstyle u}&&\downarrow\\
A&\longrightarrow&A/G .
\end{array}                                                    \tag{C.CQ.2}
\]
Here $u$ is a sequence of point blowups and $v$ is
generically finite. Pullback of a pluricanonical form
through $v$ is holomorphic. On the open set where $u$
is an isomorphism it gives a pluricanonical form on
$A$; its coefficients extend across the finitely many
missing points by Hartogs. Hence there is an injective
map
\[
 H^0(Z,mK_Z)\longrightarrow H^0(A,mK_A)
                  =\C\,\sigma^{\otimes m},
 \qquad \sigma=dz_1\wedge dz_2 .
                                                               \tag{C.CQ.3}
\]
Injectivity follows on the dense locus of the dominant
map. The image is $G$-invariant, so a nonzero image
also requires $(\det L)^m=1$.

If the quotient map $A\to A/G$ is ramified along a
curve, its generic local form is
$(x,y)\mapsto(x^e,y)$ with $e>1$. The Jacobian
pullback of a holomorphic $m$-canonical form vanishes
to order at least $m(e-1)$ along that curve. This is
incompatible with the nonzero constant form in
(C.CQ.3). We have proved
\[
 \text{a branch curve in }A\longrightarrow A/G
       \quad\Longrightarrow\quad P_m(Z)=0\quad(m>0).
                                                               \tag{C.CQ.4}
\]
This conclusion uses the actual ramification map,
including when the quotient is smooth along the
branch image.

When $L=I$, the group consists of translations. The
action is free and its quotient is the torus obtained
by enlarging $\Lambda$ by the finitely many translation
vectors. This is the first row of (C.CQ.1).

Suppose that $L$ has exactly one eigenvalue equal to
one. Every power of $L$ has that eigenvalue. A
nonidentity element of $G$ with a fixed point therefore
has a fixed curve, unless its derivative is the
identity; in the latter case it is a translation and
has no fixed point. Thus a nonfree action gives
(C.CQ.4). A free action gives a smooth finite
\'etale quotient. Its canonical line has finite order,
since the finite-order character $\det L$ is the
action on $\sigma$; it has no nonzero canonical
form. Averaging over the finite cover gives
\[
 h^1(Z,\OO_Z)=1,\qquad h^0(Z,K_Z)=0,\qquad
                  b_1(Z)=2 .
                                                               \tag{C.CQ.5}
\]
There is also an explicit product-elliptic comparison.
Let
$A_+=\ker(L-I)^0$ and $A_-=\operatorname{im}(L-I)$.
They are elliptic subtori, their tangent spaces are
complementary, and the addition map
\[
                  A_+\times A_-\longrightarrow A
                                                               \tag{C.CQ.6}
\]
is a finite \'etale isogeny. The finite group of lifts
of $G$ to this product acts freely, by translations
on the first factor and affine automorphisms on the
second; its quotient is exactly $A/G$. To see the
lift explicitly, lift the translation vector $b$
through (C.CQ.6) and use the linear action
$(1,L|_{A_-})$. Its power equal to the order of $g$
is a deck translation of (C.CQ.6), so the group
generated with the finite deck group is finite.
This proves the stated finite free product quotient,
the usual bielliptic presentation, without inferring
it only from the numerical invariants. In particular
this quotient is projective. This proves the second
row of (C.CQ.1).

It remains to treat $\ker(L-I)=0$. The endomorphism
$I-L$ of $A$ is an isogeny, because its derivative
is invertible. Solve $(I-L)p=b$ and translate the
origin to $p$. The generator becomes the linear
map $L$ and has a fixed point. Its order equals
the order of $L$.

Assume first that there is no branch curve. At the
fixed origin, no nontrivial power of $L$ can have
exactly one eigenvalue one; such a power would fix
a divisor. If the order is $n$, each eigenvalue
of $L$ is therefore a primitive $n$th root.
After changing the cyclic generator, its local
representation is
\[
             (z_1,z_2)\longmapsto
                     (\zeta z_1,\zeta^r z_2),
       \quad 1\le r<n,\quad \gcd(r,n)=1 .
                                                               \tag{C.CQ.7}
\]
We give the precise local pluricanonical test. The
quotient is the affine toric surface for the first
quadrant in
\[
 N=\Z^2+\Z(1/n,r/n),\qquad
 M=\{(a,b)\in\Z^2:a+rb\equiv0\pmod n\}.
                                                               \tag{C.CQ.8}
\]
These are the cocharacter and character lattices:
the invariant monomials are exactly
$z_1^az_2^b$ with $(a,b)\in M$. If
$(\det L)^m=1$, the form in (C.CQ.3) descends on
the smooth locus as
\[
 (z_1z_2)^m
       \left(\frac{dz_1}{z_1}\wedge
             \frac{dz_2}{z_2}\right)^{\otimes m}.
\]
Its order on the divisor of a primitive toric ray
$v=(v_1,v_2)\in N$ is
\[
                       m(v_1+v_2-1).
                                                               \tag{C.CQ.9}
\]
Indeed a unimodular toric chart expresses the
monomial's order as its character pairing with $v$,
while the logarithmic volume form has one simple
pole along that divisor. The same formula holds
on a refinement, by its explicit monomial map.

The vector $(1/n,r/n)$ is primitive in $N$ and
can be included in a smooth fan subdivision of
the quadrant. To obtain such a subdivision, first
insert this ray and then insert primitive lattice
rays until every adjacent cone has determinant
one in a lattice basis; the Euclidean algorithm
terminates this two-dimensional procedure.
When $r<n-1$, (C.CQ.9) is negative on this ray.
The descended form cannot extend to that smooth
resolution. Smooth resolutions have the same
pluricanonical sections under their actual common
point-blowup resolution, by the Hartogs comparison
already used in (C.CQ.3). Therefore
\[
 \ker(L-I)=0,\quad \det L\ne1
      \quad\Longrightarrow\quad P_m(Z)=0\quad(m>0),
                                                               \tag{C.CQ.10}
\]
also including the branch-curve case by (C.CQ.4).

Finally let $\det L=1$. Every stabilizer acts with
weights $(1,-1)$ after choosing its generator, so
there are no branch curves. For a stabilizer of
order $n$, the rays
\[
 (1,0),\ ((n-1)/n,1/n),\ldots,
                 (1/n,(n-1)/n),\ (0,1)
                                                               \tag{C.CQ.11}
\]
give a smooth subdivision. Consecutive determinants
are $1/n$, the covolume of $N$, and all these rays
have height $v_1+v_2=1$. Formula (C.CQ.9) shows
that the invariant two-form extends without a
zero or pole across every exceptional divisor.
Thus the simultaneous resolution $Z_0$ has
$K_{Z_0}\simeq\OO_{Z_0}$. Its exceptional curves
are the chains of rational curves given by
(C.CQ.11).

For clarity, its simple connectedness follows
through the quotient map as well. Remove the
discrete set of points in $\C^2$ fixed by
nontrivial elements of
$\Gamma=\Lambda\rtimes\langle L\rangle$.
The complement is simply connected: a compact
loop bounds a disc in $\R^4$, and a perturbation
relative to its boundary avoids the finitely
many removed points met in a compact region.
It is the universal cover of the smooth
quotient locus. Filling in a quotient point
kills the normal closure of its finite
stabilizer in $\Gamma$, by van Kampen.
For every $\lambda\in\Lambda$, the element
$(\lambda,L)$ has a fixed point, since $I-L$
is invertible. Thus all $(\lambda,L)$ are
killed, including $(0,L)$; their products
kill every translation $(\lambda,I)$ as well.
The resulting fundamental group is trivial.
Replacing a quotient point by the rational
chain in (C.CQ.11) imposes the same relation:
its plumbing neighbourhood is simply connected
and has the same punctured boundary as the
quotient ball. Hence $\pi_1(Z_0)=1$.
Together with $K_{Z_0}=\OO$, this identifies
$Z_0$ as a K3 surface. It is minimal, since
adjunction excludes a rational $(-1)$-curve
on a surface with trivial canonical line.
This proves the third row.

An Enriques surface has a nonzero even
pluricanonical form, but has no nonzero
canonical form and has $b_1=0$. None of
the three nonnegative cases in (C.CQ.1)
has those properties: the torus and K3
have a canonical form, and the bielliptic
case has $b_1=2$. All the other cases
have no positive pluricanonical form.
This proves the exclusion.

We finish with the component comparison.
The stabilizer of a component in the
cyclic base change of C.2.UN.16 is cyclic,
and its quotient is properly bimeromorphic
to the original normalization, as in
(C.HQ.1)--(C.HQ.3). If the minimal model
of the covering component $E$ is a torus
$A$, its actual minimal contraction
$E\to A$ is equivariant. One way to
verify the descent is to factor all the
exceptional rational curves: their maps
to the torus are constant, and the
equality $\rho_*\OO_E=\OO_A$ descends
the coordinate functions and hence
every conjugate action. Thus the induced
action on $A$ is holomorphic, and
\[
 \begin{array}{ccc}
 E&\longrightarrow&A\\
 \downarrow&&\downarrow\\
 E/G&\longrightarrow&A/G
 \end{array}
 \qquad
             E/G\dashrightarrow S_i^\nu
                                                               \tag{C.CQ.12}
\]
has a proper birational horizontal quotient
map and the specified proper birational
comparison on the right. Resolving their
graph identifies the pluricanonical
spaces and minimal birational type.
Applying (C.CQ.1) gives exactly the last
assertion. This statement concerns this
actual torus-cover branch; the separate
rank-one fibration branch is handled by
the different comparison of C.2.UN.20.
\end{proof}

For comparison with the primary literature, the projective
quotient classification is already given by
\cite[Theorem~2.1 and Lemma~3.4]{Yoshihara1995}.
Its Enriques cases have noncyclic holonomy: the Klein
four group or a dihedral group of order eight
\cite[Example~2.3 and Lemma~3.6]{Yoshihara1995}.
The proof above includes nonprojective tori and supplies
the actual cyclic component-cover comparison needed here;
no priority for the projective quotient classification
is claimed.


% C.2.UN.41 from part_C2_diagonal_hopf_occurrence.tex

